\documentclass[12pt]{ximera}
\title{Homework 8: Probability Applications of Counting}
\begin{document}
\begin{abstract}
Section 8.3: using combinations to build probabilities --- committee composition in
a state senate, and just how unlikely a magician's ``randomly shuffled'' deck really
is.
\end{abstract}
\maketitle

\section*{Section 8.3: Probability Applications of Counting Principles}

The pattern throughout this assignment: count the outcomes in the event, count all
equally likely outcomes, and divide.

\begin{problem}
A state senate currently consists of $33$ Republicans and $17$ Democrats. Its
Education committee has $15$ members --- $9$ Republicans and $6$ Democrats.

\begin{prompt}
\textbf{(a)} How many different sets of $15$ senators could possibly have been selected
for the Education committee, altogether?

$\answer{2250829575120}$

\begin{hint}
A committee is unordered, and there are $33+17=50$ senators.
\end{hint}

\begin{feedback}
Choosing an unordered committee of $15$ from $50$ senators:
\[
C(50,15)=\frac{50!}{15!\,35!}=2{,}250{,}829{,}575{,}120,
\]
over two trillion possible committees.
\end{feedback}
\end{prompt}

\begin{prompt}
\textbf{(b)} How many different ways could the $15$ have been selected so that $9$ of
them are Republicans and $6$ of them are Democrats?

$\answer{477306429600}$

\begin{hint}
Choose the Republicans and the Democrats separately, then combine the two
independent choices with the multiplication principle.
\end{hint}

\begin{feedback}
Pick $9$ of the $33$ Republicans and, independently, $6$ of the $17$ Democrats:
\[
C(33,9)\cdot C(17,6) = 38{,}567{,}100 \cdot 12{,}376 = 477{,}306{,}429{,}600.
\]
\end{feedback}
\end{prompt}

\begin{prompt}
\textbf{(c)} If the $15$ senators had been selected at random, what is the probability
that the committee would consist of $9$ Republicans and $6$ Democrats?

As a percentage rounded to one decimal place: $\answer[tolerance=0.05]{21.2}\%$

\begin{feedback}
Divide part (b) by part (a):
\[
P = \frac{C(33,9)\cdot C(17,6)}{C(50,15)}
  = \frac{477{,}306{,}429{,}600}{2{,}250{,}829{,}575{,}120}
  \approx 0.2121,
\]
about $21.2\%$.

So the actual party split of the committee is a fairly likely one --- it is close to
proportional to the Senate as a whole, and near-proportional splits are the most
common outcome of a random draw.
\end{feedback}
\end{prompt}
\end{problem}

\begin{problem}
A sleight-of-hand magician's trick involves drawing increasingly unlikely sets of
cards from a ``randomly shuffled'' deck of $52$ cards. Their deck is half red and
half black, with four suits (Hearts, Diamonds, Spades, and Clubs), and within each
suit are the cards Ace, 2, 3, 4, 5, 6, 7, 8, 9, 10, Jack, Queen, and King.

First, the magician claims, she'll draw four red cards in a row, and does. After
reshuffling the deck, she claims she'll draw four Hearts in a row, and does. After
reshuffling again, she claims she'll draw all four aces from the deck in a row, and
does. The audience is, appropriately, impressed.

\begin{prompt}
\textbf{(a)} How many distinct outcomes are possible if four cards are drawn from a
deck of $52$ and shown to the audience \emph{without} respect to their order?

$\answer{270725}$

\begin{feedback}
Order does not matter, so this is a combination:
\[
C(52,4)=\frac{52\cdot 51\cdot 50\cdot 49}{4\cdot 3\cdot 2\cdot 1}=270{,}725.
\]
\end{feedback}
\end{prompt}

\begin{prompt}
\textbf{(b)} How many distinct outcomes are possible if four cards are drawn from a
deck of $52$ and shown to the audience one at a time \emph{in the order} they were
drawn?

$\answer{6497400}$

\begin{feedback}
Now order matters, so this is a permutation:
\[
P(52,4)=52\cdot 51\cdot 50\cdot 49 = 6{,}497{,}400.
\]
This is exactly $4!=24$ times the answer to part (a), since each unordered set of
four cards can be shown in $24$ different orders.
\end{feedback}
\end{prompt}

For the remaining parts, either count can be used, as long as you count the event the
same way you counted the sample space. The solutions below use the unordered count
from part (a).

\begin{prompt}
\textbf{(c)} What is the probability, if the deck is indeed randomly shuffled, of
drawing four red cards in a row from the deck of $52$?

As a percentage rounded to one decimal place: $\answer[tolerance=0.05]{5.5}\%$

\begin{hint}
How many four-card sets consist entirely of red cards? There are $26$ red cards.
\end{hint}

\begin{feedback}
There are $26$ red cards, so the number of all-red four-card hands is
$C(26,4)=14{,}950$. Therefore
\[
P(\text{4 red}) = \frac{C(26,4)}{C(52,4)}=\frac{14{,}950}{270{,}725}=\frac{46}{833}\approx 0.0552,
\]
about $5.5\%$ --- unlikely, but hardly astonishing.
\end{feedback}
\end{prompt}

\begin{prompt}
\textbf{(d)} What is the probability, with a randomly shuffled deck, of drawing four
Hearts in a row from the deck of $52$?

As a percentage rounded to two decimal places: $\answer[tolerance=0.005]{0.26}\%$

\begin{feedback}
There are $13$ Hearts, so
\[
P(\text{4 Hearts}) = \frac{C(13,4)}{C(52,4)}=\frac{715}{270{,}725}=\frac{11}{4165}\approx 0.00264,
\]
about $0.26\%$ --- roughly $1$ chance in $379$.
\end{feedback}
\end{prompt}

\begin{prompt}
\textbf{(e)} What is the probability, with a randomly shuffled deck, of drawing all
four Aces in a row from the deck of $52$?

There is exactly one four-card hand consisting of all four aces, so this probability
is $1$ in $\answer{270725}$.

\begin{feedback}
Only one set of four cards is all four aces, namely $C(4,4)=1$, so
\[
P(\text{4 Aces}) = \frac{1}{C(52,4)}=\frac{1}{270{,}725}\approx 0.0000037,
\]
about $0.00037\%$.

Taken together, the three tricks are the point of the problem: the last one alone is
a one-in-$270{,}725$ event, so the deck was almost certainly not randomly shuffled.
\end{feedback}
\end{prompt}
\end{problem}

\end{document}
